\chapter{Implementation}\label{ch:implementation}

This chapter defines all implementation choices that were made to validate the research question and test the hypotheses defined in \ref{sec:research-question}.

\section{Tech stack}
\subsection{Hardware}
\subsubsection{Computer}
The development machine was a MacOS laptop with 32 GB of memory and 994 GB of storage. The storage was important to install and run the WebArena configurations, since they ship as Docker images of four full websites and take up a lot of space. This machine ran the private and shared memory generation, against Apple's MLX runtime and the Metal Performance Shaders (MPS) backend. Both models are used here: Gemma for the differentially private token-by-token generation, and Qwen for the title and description that follow it.

\subsubsection{GPUs}
The GPUs that are available to us through the school are a cluster of three NVIDIA A30s. These were used to run the agent, and therefore the Think node, across the 812 WebArena tasks. Although this was not the quickest processing unit available, it was all that was available on a free basis, and each experimental condition requires a complete pass-through of all 812 tasks, so the throughput gain over the laptop compounds across runs. These GPUs were shared across multiple students, and therefore had to be used efficiently.

The work is therefore split across the two machines by what each part needs. The agent loop is throughput-bound, since it repeats across 812 tasks and once per experimental condition, and it runs on the cluster. The differentially private path is bound instead by the token-by-token logit vector prediction over multiple memories, which is very time-consuming, and it runs on the laptop, where the model can be held in memory across a whole generation run without competing for shared cluster time. This split influenced a number of choices on the software side.

\subsection{Software}

\subsubsection{Programming language}
The programming language for this entire project is Python. The reason for choosing it was for multiple reasons: first, it is the programming language the author is most well-versed in. Second, Python hosts a mature scientific stack, built on the array-programming foundation of NumPy \citep{harris2020array} and including the deep learning framework this thesis depends on for local logit access, PyTorch \citep{paszke2019pytorch}. Both the differentially private mechanism and the agent loop are therefore assembled from maintained, interoperable libraries rather than reimplemented.

\subsubsection{Embedding model}
The thesis embeds text with \texttt{all-MiniLM-L6-v2}, a sentence-transformers model \citep{reimers2019sentencebert} built as a six-layer reduction of the distilled MiniLM architecture \citep{wang2020minilm} and fine-tuned with a contrastive objective on roughly 1.17 billion sentence pairs. It is used in two places: to retrieve memories for the Think prompt, and to assign each memory to one of the differentially private labels in the shared memory pipeline described in Section \ref{sec:differential privacy mechanism}.

The model was chosen for its cost rather than its ceiling. At 22.7 million parameters and 384 dimensions it runs offline and deterministically on the laptop, and the retrieval step is called once per Think turn across 812 tasks per pass-through, so the per-call cost compounds quickly. MTEB \citep{muennighoff2023mteb} reports that no single embedding model dominates across task types, so the choice is a trade-off rather than a ranking, and \texttt{all-MiniLM-L6-v2} sits at the small and fast end of it. While a larger encoder would likely retrieve more accurately, this thesis compares configurations against each other under one fixed retrieval setting, so holding retrieval quality constant across conditions matters more than maximising it.

\subsubsection{LLM models}
The LLM models that this thesis uses were chosen with the aim of firstly, meeting the requirements for the use in terms of output, secondly, of being free of charge, thirdly, of producing high-quality output, and fourth, of producing output in a reasonable time-frame. Lastly, this thesis gave preference to models used by either the ReasoningBank or Amin et al. approach, since the aim is to replicate their mechanisms as closely as possible. These objectives had to be traded off against each other for the different use cases.

LLMs were used in the thesis in three parts:
\begin{enumerate}
    \item \textbf{Agent - Thinking node.} The thinking node in the ReAct loop
    requires an LLM to take the information from the past actions and
    observations and make a choice on the next action from a subset of
    pre-defined actions. This step does not require any intermediate
    computations, and simply requires a complete answer: a single line naming
    one action call. The choice of model could therefore be focused on a free
    and quick model.

    In a first step, the author checked for the model that was used to run the
    agent in the ReasoningBank paper. Unfortunately, the paper used
    closed-source, paid models: Gemini-2.5 and Claude 3.7. Since there was no
    budget for this project, and any low-tier free access would have been
    exhausted well before 812 tasks completed, this thesis had to take a
    different approach\footnote{This is an important distinction to the tests
    run by the ReasoningBank paper \citep{ouyang2025reasoningbank}. It would
    therefore not be unexpected if the model underperformed in comparison to
    that paper on the 812 benchmarks.}.

    For open-source models, Qwen is among the most widely adopted open-weight
    model families, with its technical reports accumulating several thousand
    citations \citep{yang2024qwen25, yang2025qwen3} and its checkpoints
    released under Apache 2.0. It could be downloaded for free, and had
    different parameter settings to choose from to balance speed and
    intelligence. The model used for the Think node is Qwen3.5 at 4 billion
    parameters, quantised to NVFP4 (Ollama tag
    \texttt{qwen3.5:4b-nvfp4}).\footnote{Qwen3.5 was released without a
    technical report; the architecture and training lineage are documented for
    the preceding generation in \citet{yang2025qwen3}.}

    The reason for choosing the 4bn parameter model over the 9bn one, is the
    speed of execution. The 9bn parameter model was taking overall too long,
    whereas the 4bn parameter model allowed an average of $\sim$5\,s per Think
    step. Two limits on that figure should be stated. It is a wall-clock
    measurement of one configuration rather than evidence about any single
    factor, because moving to it changed the parameter count and the
    quantisation format at the same time, and no controlled ablation was run.
    It was also measured on the laptop during model selection, whereas the
    evaluation runs were served from the A30 cluster, so it is the basis on
    which the 4bn model was chosen and not the per-step cost of the reported
    runs.

    Since the thesis aim was not to assess overall success rates, but a
    comparative assessment instead, the overall intelligence of the model was
    not as important. Published zero-shot success rates for open-weight models
    at or below 8 billion parameters on the full 812-task benchmark span
    roughly 2\% to 24\%, and that spread is driven at least as much by the
    agent harness, the observation format and the prompting strategy as by the
    model itself \citep{zhang2025symbiotic, gandhi2026gobrowse, lu2026distillation}.
    The closest published reference point for the configuration used here is
    \citet{lu2026distillation}, who report the same Qwen3.5 4bn model at 24.1\%
    on the full 812 tasks under the BrowserGym evaluation protocol, without
    fine-tuning. That figure is a useful target rather than a prediction, since
    the harness differs from the one built for this thesis. While a low
    absolute success rate would narrow the base of tasks on which a difference
    between configurations can show, it would not undermine the comparison,
    which holds the model fixed and varies only the memory condition.

    The model is served locally through an Ollama endpoint. Ollama runs the
    quantised weights as a local HTTP service, so each runner calls one chat
    endpoint rather than loading its own copy of the model, and no request
    leaves the machine. Its chat endpoint returns the generated text and
    nothing else, which becomes the deciding constraint for the differentially
    private path below.

    \item \textbf{Differential privacy mechanism.} The model used to generate
    the shared memories could no longer be the Qwen model. The Amin et al.
    mechanism needs the full next-token logit vector over the vocabulary at
    every sampling step, for each member of the batch and for the public
    prompt, and it needs to be able to extend that generation one token at a
    time. Qwen is served in this thesis through Ollama's chat endpoint, which
    returns the generated text and nothing else, so the logit vectors are never
    exposed. The differentially private path consequently runs on a locally
    loaded Gemma model, through PyTorch and HuggingFace transformers, where the
    logits are read directly off the forward pass.

    Naively, each sampled token would cost one full forward pass over the
    prompt for every sensitive prompt plus one for the public prompt. The
    implementation instead stacks all of them into a single padded batch, runs
    one prefill with the key-value cache retained, and extends the generation
    one token at a time through a single-column continuation that reuses that
    cache, so the unchanging prompt prefix is attended to once rather than at
    every step.

    The sole deviation of substance from the \citet{amin2024private} paper is
    that it uses Gemma 1.1 2B IT, whereas this thesis uses the more recent
    Gemma 2 2B IT \citep{gemmateam2024gemma2}. This choice may affect the
    hyperparameters, since the values the paper reports were tuned on the
    earlier model, and the values used here are given in Table
    \ref{tab:dp-hyperparameters} below.

    Gemma 2 2B IT was chosen over the Gemma 2 2B base model for two reasons.
    First, it has been instruction-tuned, meaning further trained on tasks
    phrased as natural-language instructions \citep{wei2022finetuned}, so
    inputs must be arranged into the model-specific chat template. Second, in
    the intermediate experiments reported in Section \ref{intermediate results},
    Gemma 2 2B IT under the chat template produced lessons that fitted the
    input memory items better than the base model under raw completion. That
    comparison was a single qualitative pass over a ten-item fixture, with the
    model and the prompting regime varied together, so it is reported as the
    reason the project standardised on the IT variant rather than as a measured
    result.

    The model is loaded in \texttt{torch.bfloat16} rather than fp32 or fp16.
    Gemma 2 applies logit soft-capping in each attention layer and in the final
    layer \citep{gemmateam2024gemma2}, and fp16's narrower exponent range
    produces silent NaN and Inf values there, while fp32 was too slow on MPS.
    bfloat16 is the model's training dtype and preserves the fp32 dynamic range
    at roughly half the per-forward cost.

    \item \textbf{Shared memory title and description generation.} The Qwen
    model is used once more, for the title and description generation of the
    shared memory. This step no longer requires token logit access, so it can
    be served through the same Ollama endpoint as the Think node. Its inputs
    are strictly the differentially private content and the differentially
    private label, so by the post-processing property of differential privacy
    the call adds nothing to the privacy account.

    The choice of Qwen over Gemma at this step is driven by role rather than by
    a measured capability gap. Gemma is loaded for the differentially private
    path because that path needs token-level logit access, which is a
    requirement of the mechanism and not a judgement about the model. Once the
    content has been released, the remaining work is short-form instruction
    following, and Qwen is already resident in the pipeline behind the same
    \texttt{ChatClient} interface the Think node uses, so reusing it costs
    nothing. No head-to-head benchmark between the two is reported here: a 4
    billion parameter model released in 2026 and a 2 billion parameter model
    released in mid-2024 differ in scale, training recipe and release date at
    once, so a comparison between them would be confounded on three axes and
    its outcome is not in question.
\end{enumerate}

\subsubsection{Testing configuration}
The testing configuration consisted of both the tasks and scoring, as well as the environment. The tasks and scoring were provided by WebArena, while the environment was set by BrowserGym.

\textbf{WebArena}

All evaluation runs are completed on the WebArena gym \citep{zhou2024webarena},
accessed through the BrowserGym wrapper \citep{dechezelles2024browsergym}. The
WebArena repository provides the evaluation environment, tasks, and scoring for
812 different online tasks:

\begin{itemize}
    \item \textbf{Environment.} The WebArena repository provides the
    configurations for four websites, which mimic the real-world equivalent of
    GitLab, Reddit, an online store content management system (CMS), and
    OneStopShop. These configurations allow the tester to run the agent tests
    locally on replicas of the websites, removing the risk that the agent's
    actions take effect on a live service. It also provides a set of utility
    tools and knowledge resources that the agent can call, if needed, for
    further closeness to the real world scenario. These include a calculator, a
    scratchpad, an offline copy of Wikipedia, and a map, the last of which is
    itself served as a site and carries its own tasks \citep{zhou2024webarena}.
    Because the environment resets to a deterministic initial state, these
    configurations also ensure that the same tests can be run over an extended
    duration of time, without fear of change to the website.

    \item \textbf{Tasks.} The repository also provides a list of 812 different
    online tasks, which the agent is given to perform. These tasks are
    instantiated from 241 templates and are phrased as high-level, realistic
    natural language commands rather than click-by-click instructions, so the
    agent has to decompose them itself \citep{zhou2024webarena}. The example
    task \textit{``What is the top-1 best-selling brand in Quarter 1 2022''}
    names an outcome and leaves the route to it open.

    \item \textbf{Evaluation.} Lastly, the repository also provides a program,
    which validates the functional correctness of each task by checking the
    outcome rather than the action sequence taken to reach it. BrowserGym
    surfaces that verdict as the episode reward, and the runner records it when
    the episode ends, so a task counts as successful when its final reward is
    greater than zero.
\end{itemize}

The reason for choosing the WebArena environment was to keep as close as possible to the original testing environment used by \citet{ouyang2025reasoningbank}. The thesis therefore chose to test its hypotheses on the same test set as that paper. Firstly, this ensures that the ReasoningBank framework is useful on this dataset, and that a difference in agent behaviour with memories should be observable. Second, it allows the results of this thesis to be compared to those in the paper.

\textbf{BrowserGym}

While WebArena provides the benchmark tasks and evaluations, BrowserGym provides the testing environment to run the thesis on. It launches and takes action on a Chromium browser using Playwright. While the work for this thesis included creating its own Playwright nodes and Chromium actions, BrowserGym simplifies the work in three ways. First, it exposes the WebArena suite through a single environment call, including the per-task setup and the reward signal that the hand-rolled path had no access to. Second, it offers a larger action vocabulary: the WebArena subset used here adds hovering, dropdown selection, multi-tab handling, history navigation, and a structured channel for returning an answer or reporting a task infeasible, against the seven actions the hand-rolled implementation supported. Third, it catches exceptions raised by an action and returns the error text in the next observation \citep{dechezelles2024browsergym}, so the failure is fed back to the LLM rather than lost.

It also returns structured observations, in one of three formats. BrowserGym automatically returns all three observation spaces which WebArena has been tailored to \citep{zhou2024webarena, dechezelles2024browsergym}: the HTML DOM tree, a screenshot, or the accessibility tree of the website URL and content. Which of the three this thesis feeds into the Think node, and the reason for that choice, is set out in Section \ref{sec:experiment_design}.
